% 読み物版の共通プリアンブル。各章の .tex から \input する。
\usepackage{amsmath}
\usepackage{bm}
\usepackage{array}
\usepackage{booktabs}
\usepackage{colortbl}
\usepackage{multirow}
\usepackage{graphicx}
\graphicspath{{fig/}{../fig/}}   % 公開用のソースは book/fig/、正本は ../fig/ に図がある
\usetikzlibrary{arrows.meta,positioning}
\tikzset{
  kbox/.style   = {draw=kkaccent, fill=kkpale, line width=.5pt,
                   align=flush center, inner sep=5pt, font=\footnotesize,
                   rounded corners=1.5pt, minimum height=8mm},
  kplain/.style = {kbox, fill=white},
  kar/.style    = {-{Stealth[length=4.5pt,width=3.5pt]}, draw=kkaccent, line width=.6pt},
  knode/.style  = {circle, draw=kkaccent, fill=kkpale, line width=.6pt,
                   minimum size=8mm, inner sep=0pt, font=\small},
  kedge/.style  = {draw=kksecondary, line width=1pt},
}

% jlreq は行送りを固定するため、背の高い箱（表・図）が前後の行と重なる。
% 表と図はこの環境で囲み、前後に必ずアキを入れる。
\newenvironment{kkcenter}
  {\par\addvspace{\baselineskip}\begin{center}}
  {\end{center}\par\addvspace{.5\baselineskip}}


% 版面幅を超える表は自動で縮める。収まっていれば原寸のまま。
\newsavebox{\kkfitbox}
\newcommand{\kkfit}[1]{%
  \sbox\kkfitbox{#1}%
  \ifdim\wd\kkfitbox>\linewidth
    \resizebox{\linewidth}{!}{\usebox\kkfitbox}%
  \else
    \usebox\kkfitbox
  \fi}

\newcommand{\kkcap}[1]{\\[.5\baselineskip]{\footnotesize\kkgt{3}\textcolor{kkdark}{#1}}}

% 読み物版では節見出しの前アキを詰める（既定は 2 行ドリ）
\makeatletter
\renewcommand{\kk@heading@pre}{%
  \ifkk@lastisheading \vspace{-.3\baselineskip}\else \vspace{1\baselineskip}\fi}
\makeatother
\raggedbottom

\usepackage{hyperref}
\hypersetup{luatex, pdfencoding=auto, colorlinks=true,
            linkcolor=black, urlcolor=blue}

% 章の扉。合本にしたときの区切りに使う。
\makeatletter
\newcommand{\kkchapter}[2]{%
  \clearpage
  \null\vspace{6mm}%
  \noindent{\color{kkaccent}\rule{\linewidth}{1.4pt}}\par\vspace{4mm}%
  \gdef\kkcurchap{#1}%
  \noindent{\kkgt{4}\large\kkchaplabel{#1}}\par\vspace{3mm}%
  \noindent{\kkgt{7}\fontsize{21pt}{27pt}\selectfont #2}\par\vspace{4mm}%
  \noindent{\color{kksecondary}\rule{\linewidth}{.5pt}}\par
  \vspace{6mm}%
  \setcounter{section}{0}%
  \global\kk@lastisheadingfalse
}
\makeatother

% ---- メタ解析入門で追加 ----------------------------------------------------
\usepackage{fvextra}
\usepackage{pgfplots}
\pgfplotsset{compat=1.18}
\usetikzlibrary{shapes.geometric,calc,fit,backgrounds,patterns,decorations.pathreplacing}
\definecolor{kkwarm}{HTML}{C0504D}
\definecolor{kkcode}{HTML}{F4F7FA}

% R のコード。そのまま写して実行できる形で載せる
\DefineVerbatimEnvironment{rcode}{Verbatim}{%
  breaklines, breakanywhere, fontsize=\small, frame=leftline, framerule=1.2pt,
  rulecolor=\color{kkaccent}, framesep=3mm, xleftmargin=1mm,
  bgcolor=kkcode}
% R の出力。コードと見分けがつくよう地の色を外し、小さく組む
\DefineVerbatimEnvironment{rout}{Verbatim}{%
  breaklines, breakanywhere, fontsize=\footnotesize, frame=single, framerule=.4pt,
  rulecolor=\color{kksecondary}, framesep=2mm, formatcom=\color{kkdark}}

% 本文の TikZ 図は、紙面より広いときだけ縮める（Kindle 版の狭い紙面用。A5、A4 では何も変わらない）
\usepackage{adjustbox}
\newenvironment{kkfitpic}{\par\noindent\begin{adjustbox}{max width=\linewidth,center}}{\end{adjustbox}}

% 図を入れる（kkcenter と見出しを一緒に）
% R の図は浮動体にする。次のページへ送られても、空いた所に本文が流れる
\newcommand{\kkfig}[3][\linewidth]{%
  \begin{figure}[!htbp]\centering
  \includegraphics[width=#1]{#2}\kkcap{#3}\end{figure}}
\renewcommand{\topfraction}{.85}\renewcommand{\bottomfraction}{.7}
\renewcommand{\textfraction}{.12}\renewcommand{\floatpagefraction}{.75}
\setlength{\textfloatsep}{12pt plus 4pt minus 2pt}

% 部の扉
\newcommand{\kkcurchap}{}
\newif\ifkkappendix
\newcommand{\kkchaplabel}[1]{\ifkkappendix 付録 #1\else 第 #1 章\fi}
\newcommand{\kkpart}[3]{%
  \clearpage\thispagestyle{empty}%
  \null\vspace{40mm}%
  \if\relax\detokenize{#1}\relax\else\noindent{\kkgt{4}\large\color{kkaccent} 第 #1 部}\par\vspace{4mm}\fi
  \noindent{\kkgt{7}\fontsize{26pt}{32pt}\selectfont #2}\par\vspace{6mm}%
  \noindent{\color{kkaccent}\rule{\linewidth}{1.4pt}}\par\vspace{8mm}%
  \noindent\begin{minipage}{\linewidth}\small #3\end{minipage}%
  \clearpage}

% 章の頭の「この章で分かること」
\newenvironment{chapintro}{\begin{simpleao}{この章で分かること}}{\end{simpleao}}

% 節と小節の見出しを箱なしで組み直す。tcolorbox の見出しは空きがあっても次ページへ送られるため
\makeatletter
\renewcommand{\kk@section@typeset}[2]{%
  \par\noindent{\fontsize{11pt}{14pt}\selectfont\kkgt{5}\color{kkaccent}#1}\par\nobreak
  \vspace{.5mm}%
  \noindent\parbox{\linewidth}{\raggedright\fontsize{15pt}{20pt}\selectfont\kkgt{4}#2}\par\nobreak
  \vspace{1.2mm}%
  \noindent{\color{kkaccent}\rule{\linewidth}{.8pt}}\par\nobreak
  \vspace{2.5mm}\@afterheading}
\renewcommand{\kk@subsection@typeset}[2]{%
  \par\noindent{\color{kkaccent}\rule[.15ex]{1.2mm}{1.2mm}}\hspace{2mm}%
  {\fontsize{12pt}{16pt}\selectfont\kkgt{4}#2}\par\nobreak
  \vspace{1.5mm}\@afterheading}
\makeatother

% \raggedbottom では負のペナルティの位置が常に最良の改ページ点になる。
% リストの前後で改ページされて下に空白が残るのを防ぐため、リストのペナルティを 0 にする
\makeatletter
\@beginparpenalty=0 \@endparpenalty=0 \@itempenalty=0
\makeatother

% 初読では飛ばしてよい内容の印。節には見出しの頭に \hatten を置き、段落には hattenbox 環境を使う
\DeclareRobustCommand{\hattentag}{\tikz[baseline=(t.base)]\node[draw=kkaccent, rounded corners=1pt,
  inner xsep=1.6pt, inner ysep=1.2pt, text=kkaccent, font=\scriptsize\kkgt{5}](t){発展};}
\newcommand{\hatten}{\texorpdfstring{\hattentag\hspace{2mm}}{[発展] }}
\newtcolorbox{hattenbox}{breakable, enhanced, colback=white, colframe=kkaccent!45,
  boxrule=.4pt, arc=1pt, left=2.5mm, right=2.5mm, top=1.5mm, bottom=1.5mm,
  before skip=3mm, after skip=3mm, before upper={\hattentag\hspace{2mm}}}

% 目次に部と章の行を足す。kkchapter と kkpart は見出しを自前で組むため、目次の行をここで書く
\makeatletter
\newcommand{\l@kkpartline}[2]{\addvspace{7pt}\noindent{\kkgt{5}\small\color{kkaccent}#1}\par\nobreak}
\newcommand{\l@kkchapline}[2]{\addvspace{4pt}\noindent
  {\kkgt{5}\small\color{kkdark}#1}\kk@toc@leader{\small\color{black!75}#2}\par\nobreak}
\let\kk@orig@chapter\kkchapter
\renewcommand{\kkchapter}[2]{\kk@orig@chapter{#1}{#2}%
  \addcontentsline{toc}{kkchapline}{\kkchaplabel{#1}\hspace{.6\zw}#2}}
\let\kk@orig@part\kkpart
\renewcommand{\kkpart}[3]{%
  \clearpage
  \addcontentsline{toc}{kkpartline}{\if\relax\detokenize{#1}\relax\else 第 #1 部\hspace{.6\zw}\fi #2}%
  \kk@orig@part{#1}{#2}{#3}}
\let\kk@orig@toc\tableofcontents
\renewcommand{\tableofcontents}{%
  \thispagestyle{empty}%
  \noindent{\kkgt{7}\fontsize{17pt}{21pt}\selectfont 目次}\par\vspace{1.5mm}%
  \noindent{\color{kkaccent}\rule{\linewidth}{1pt}}\par\vspace{2mm}%
  {\footnotesize 見出しの前の\hattentag の印は、初めて読むときは飛ばしてよい部分です。}\par\vspace{3mm}%
  \kk@orig@toc}
\makeatother

% 用語集で、見出し語の横に読み方を小さな字で添える
\newcommand{\kkyomi}[1]{\hspace{.5\zw}{\footnotesize\mdseries\textcolor{kkdark}{（#1）}}}

% 相互参照。節の番号に章の番号を含め、「5.2 節（p. 42）」の形で参照できるようにする
% 節と小節には \label{sec:...} を付け、本文では \kksee{sec:...} か \kkref{sec:...} で参照する
% 章には \kkchapter が ch:5、ch:A のようなラベルを自動で付ける。参照は \kkchref{5}
\makeatletter
\renewcommand*\thesection{\kkcurchap.\arabic{section}}
\let\kk@orig@chapter@b\kkchapter
\renewcommand{\kkchapter}[2]{\kk@orig@chapter@b{#1}{#2}%
  \protected@edef\@currentlabel{#1}\label{ch:#1}}
\makeatother
\newcommand{\kkref}[1]{\ref{#1} 節（p.~\pageref{#1}）\ignorespaces}
\newcommand{\kksee}[1]{（→ \ref{#1} 節（p.~\pageref{#1}））\ignorespaces}
\newcommand{\kkchref}[1]{\kkchaplabel{#1}（p.~\pageref{ch:#1}）\ignorespaces}
% hyperref のリンク先を章ごとに区別する（既定では各章の 1 節が同じ識別子になり、別の章へ飛ぶ）
\renewcommand*\theHsection{\kkcurchap.\arabic{section}}
\renewcommand*\theHsubsection{\kkcurchap.\arabic{section}.\arabic{subsection}}
